\scanpage{120}
\mspage{119}


``compatible with èommutative $\unsure{\otimes|\Omega}$'' is a functor and a $c$, compatible with the datum of the$\unsure{\otimes|\deleted{x}}$ $\handwritten{\unsure{\mu|\psi}}$ .

When one has a law $\otimes$ with data a) and b), it is appropriate to postulate a compatibility between the two (in the framework of axiomatics 1), is this a consequence of the other axioms ? I \deleted{fear}\inserted{\handwritten{think}} not, and that it must be added). When one has a law $\otimes$ with data a) and \rlap{v}c), it is appropriate to postulate a compatibility between the two. Likewise for a law $\otimes$ with b) and c). When one has a)b)c) at the same time, the package \deleted{\unsure{of axioms}} of all the compatibility conditions are those stated above \inserted{in 1º\handwritten{\begin{tikzpicture}[x=.45cm,y=.45cm,baseline=-.2ex]\draw (0,-.3) .. controls (-.12,.14) and (.05,.3) .. (.29,.28)--(1.5,.32) .. controls (1.71,.35) and (1.77,.53) .. (1.73,.66);\end{tikzpicture}}} a) to d) . Likewise, a Functor``compatible with associative $\otimes$ with two-sided \unsure{\rlap{\begin{tikzpicture}[x=1em,y=1em,baseline=0pt]\draw[line width=.45pt] (.38,.04)--(.38,.45) .. controls (.43,.42) and (.48,.29) .. (.42,.24);\fill (.28,.72) circle (.045);\fill (.82,.77) circle (.03);\fill (.96,.75) circle (.04);\draw[line width=.4pt] (.7,.08)--(.62,.27)--(.73,.43);\end{tikzpicture}}unit|moth-eaten}'' is a functor $F$, $+$ a $c$, compatible separately with associativity and unit two-sided\rlap{e}ness, and likewise for any other package of additional data among a),b),c). Finally, the notion of a homomorphism of \deleted{functo} $\otimes$-functors does not depend on the additional structures a)\deleted{b)} or c) that one can consider\inserted{\handwritten{(but indeed on b)\begin{tikzpicture}[x=.6cm,y=.6cm,baseline=-.2ex]\draw (0,-.25) .. controls (.03,.05) and (.11,.15) .. (.31,.17)--(2.6,.19) .. controls (2.86,.21) and (3.05,.32) .. (3.04,.56);\end{tikzpicture}}} , being expressed exclusively by means of the datum of $c$\deleted{\unsure{)}}\handwritten{\unsure{\begin{tikzpicture}[x=.45cm,y=.45cm,baseline=-.3ex]\draw[line width=.8pt] (.48,.9) .. controls (.22,.68) and (.25,.64) .. (.55,.66) .. controls (.72,.68) and (.68,.49) .. (.61,.45);\draw[line width=.8pt] (.38,.61) .. controls (-.03,.54) and (.02,.39) .. (.15,.29)--(.55,.29);\draw[line width=1pt] (.16,.3) .. controls (-.04,.1) and (.02,.01) .. (.18,-.09)--(.16,.3);\end{tikzpicture}}} \deleted{(subject t\unsure{\rlap{\raisebox{1ex}{\hspace{.25em}\begin{tikzpicture}[x=1em,y=1em]\draw[line width=1pt] (0,0)--(0,.28);\end{tikzpicture}}}o|o} verification indicated in 3º)}. - \deleted{When}\handwritten{As} for the \deleted{bilinear structure, it is appropriate to}\inserted{\deleted{\unsure{biaddii|biaddiv}}} condition of additivity on the categories and of biadditivity on $\otimes$, it is appropriate to introduce it only at the very end.

\margin{Careful whether associativity ought not to come back in here too}
\deleted{5)}\inserted{\handwritten{5)}} $\mathbb{NB}$ If one has $\otimes$ with two-sided unit, then for two endomorphisms $u,v$ of $\underline{1}$, one has $u\otimes v\overset{\handwritten{!}}{=}uv=vu$. Thus under the commutativity condition\deleted{\unsure{x|,}} c), $\unsure{\rlap{\(\mathrm{e}\)}\mathrm{E}|\mathbb{E}}\mathrm{nd}(\underline{1})$ is a commutative monoid. \inserted{\deleted{it acts on all}}\deleted{If $\underline{C}$ is additive, it is} Without a commutativity condition this monoid acts on the left and on the right on every object of $\underline{C}$; the two actions co\rlap{n}incide assuming commutativity. If \deleted{the structure}\inserted{$C$} is add\rlap{o}itive, $\mathrm{E}\rlap{\(\mathrm{b}\)}\mathrm{nd}(\underline{1})$ is a ring, which acts on the left and on the right on each object of $C$. In the commutative case, $\mathrm{End}(\underline{1})=\unsure{\mathbb{k}|k}$ is therefore a comm\rlap{u}ytative ring and the category $C$ is $k$-linear in a natural way\handwritten{, $\otimes$ is $k$-bilinear.}
